%----------------------------------------------------------------------------------------------
%------------------------------     1. Lin Gleichungssysteme    -------------------------------
%----------------------------------------------------------------------------------------------
\section{Lineare Gleichungssysteme}


%----------------------------------------------------------------------------------------------
%------------------------------     1.1 Matritzen und Vektoren    -----------------------------
%----------------------------------------------------------------------------------------------
\subsection{Matritzen und Vektoren}
\begin{tabular}{ll ll}
	
	Matrix: &  
		$A = 
		\begin{pmatrix}
			a_{11}&a_{12}&a_{13}\\
			a_{21}&a_{22}&a_{23}\\
			a_{31}&a_{32}&a_{33}
		\end{pmatrix}$ 
		& 
	Einheitsmatrix: &
	 $I = E = 
		\begin{pmatrix}
			1&0&0\\
			0&1&0\\
			0&0&1
		\end{pmatrix}$ 
	\\
	Zeilenvektor: &
	$a = 
	\begin{pmatrix}
		a_1&a_2&a_n
	\end{pmatrix}$ 
	& 
	Spaltenvektor:&
	$Ax = b = 
	\begin{pmatrix} 
		b_1\\b_2\\b_3 
	\end{pmatrix}$ 
\end{tabular}

\begin{tabular}{l lll}
	Standardbasisvektoren: &
	$ e_1= e_x =
	\begin{pmatrix}	
		1\\0\\0
	\end{pmatrix},$ &
	$ e_2= e_y =
	\begin{pmatrix}
		0\\1\\0
	\end{pmatrix},$ &
	$ e_3= e_z =
	\begin{pmatrix}
		0\\0\\1
	\end{pmatrix}$
\end{tabular}


%----------------------------------------------------------------------------------------------
%---------------------     1.2 Rechenregeln mit Matritzen und Vektoren    ---------------------
%----------------------------------------------------------------------------------------------
\subsection{Rechenregeln mit Matritzen und Vektoren}

\subsubsection{Addition / Subtraktion}

\begin{tabular}{ll}
	Beschreibung für Spaltenvektoren 	& \textrightarrow\ gilt analog für Zeilenvektoren!	\\
	Beschreibung nur für Addition 		& \textrightarrow\ gilt analog für Subtraktion!
\end{tabular}

\vspace{0.2cm}

\begin{tabular}{lll}
	$a = 
	\begin{pmatrix} 
		1\\2\\3 
	\end{pmatrix}$ & 
	$b = 
	\begin{pmatrix} 
		4\\5\\6
	\end{pmatrix}$ &
	$a+b = 
	\begin{pmatrix} 
		1+4\\2+5\\3+6 
	\end{pmatrix}$ 
	\\
	\\
	$A=
	\begin{pmatrix}
		1&2\\
		3&4\\
		5&6
	\end{pmatrix}$ &
	$ B=
	\begin{pmatrix}
		0&7\\
		5&3\\
		6&2
	\end{pmatrix}$ &
	$A+B=
	\begin{pmatrix}
		1+0&2+7\\
		3+5&4+3\\
		5+6&6+2
	\end{pmatrix}$
\end{tabular}


		    
\subsubsection{Multiplikation}

%Vektor mal Faktor
\begin{tabular}{ll}
	$v = 
	\begin{pmatrix}
		1\\8\\4 
	\end{pmatrix}$ &
	$\lambda \cdot v = 
	\begin{pmatrix} 
		\lambda \cdot 1\\
		\lambda \cdot 8\\
		\lambda \cdot 4
	\end{pmatrix}$ 
\end{tabular}

\vspace{0.2cm}

%Matrix mal Faktor 
\begin{tabular}{ll}
	$D =
	\begin{pmatrix}
		6&1&2\\
		5&4&0\\
		4&1&9
	\end{pmatrix}$ &
	$\lambda \cdot D =
	\begin{pmatrix}
		\lambda \cdot 6 & \lambda \cdot 1 & \lambda \cdot 2 \\
		\lambda \cdot 5 & \lambda \cdot 4 & \lambda \cdot 0 \\
		\lambda \cdot 4 & \lambda \cdot 1 & \lambda \cdot 9
	\end{pmatrix}$ 
\end{tabular}

\vspace{0.2cm}	

% Matrix mal Vektor			
$D \cdot v = 
	\begin{pmatrix}
		\color{red}6&\color{red}1&\color{red}2\\
		5&4&0\\
		4&1&9
	\end{pmatrix}
	\begin{pmatrix}
		\color{blue}1\\
		\color{blue}8\\
		\color{blue}4
	\end{pmatrix}
	=
	\begin{pmatrix}
		{\color{red}6}\cdot{\color{blue}1}+
		{\color{red}1}\cdot{\color{blue}8}+
		{\color{red}2}\cdot{\color{blue}4}
		\\
		5\cdot{\color{blue}1}+
		4\cdot{\color{blue}8}+
		0\cdot{\color{blue}4}
		\\
		4\cdot{\color{blue}1}+
		1\cdot{\color{blue}8}+
		9\cdot{\color{blue}4}
	\end{pmatrix}
	=
	\begin{pmatrix}
		22\\37\\48
	\end{pmatrix}$ 

	\vspace{0.2cm}

% Matrix mal Matrix
\begin{tabular}{ll}
	$A =
	\begin{pmatrix}
		\color{red}1&\color{red}2&\color{red}4\\
		0&4&-1
	\end{pmatrix}$, &
	$B =
	\begin{pmatrix}
		\color{blue}-2&4\\
		\color{blue}-4&-5\\
		\color{blue}2&4
	\end{pmatrix} $
\end{tabular}
	
$A \cdot B =
	\begin{pmatrix}
		{\color{red}1}\cdot({\color{blue}-2})+{\color{red}2}\cdot({\color{blue}-4})+{\color{red}4}\cdot{\color{blue}2} &
		{\color{red}1}\cdot 4+{\color{red}2}\cdot(-5)+{\color{red}4}\cdot 4\\
		0\cdot ({\color{blue}-2})+4\cdot({\color{blue}-4})+(-1)\cdot {\color{blue}2} &
		0\cdot 4+4\cdot(-5)+(-1)\cdot 4
	\end{pmatrix}$


\subsubsection{Division}
			
Divisionen von Matritzen werden durch eine \textbf{Multiplikation mit der inversen Matrix} durchgeführt.


\subsubsection{Zusammenfassung / Erweiterung der Rechenregeln}
   			
\begin{tabular}{lll}
	Symmetrie: 			& $ A = A^t$		 												\\   			
	Orthagonalität: 	&  $ A A^t = A^t A = I $  			& $A^{-1} = A^t$				\\
						& $A^t$ auch orthagonal 			& $\det(A) = \pm 1$ 			\\
	Produkte: 			& \textbf{AB $\neq$ BA} 			& $(AB)C = A(BC)$				\\
						& $IA = A = AI$  					& $A(B + C) = AB + AC$			\\
						& $A (\lambda B) = (\lambda A) B$ 							\\
	Transponiert: 		& $(AB)^t = B^t A^t $ 												\\
	Inverse: 			& $A^{-1} A = I$ 					& $(AB)^{-1} = B^{-1} A^{-1} $ 	\\
	Potenzen: 			& $(A^t)^{-1} = (A^{-1})^t $
\end{tabular}
